% BHU PYQ -- Transcribed & Corrected
% Paper: CHESE11: Quality Control in Chemical Analysis
% Exam: B.Sc. (Hons) Semester I (NEP) Examination 2024-25
% Category: Skill Enhancement Course (SEC)
% Department: Chemistry

\documentclass[11pt,a4paper]{article}
\usepackage[a4paper,top=2.5cm,bottom=2.5cm,left=2.8cm,right=2.8cm,headheight=14pt]{geometry}
\usepackage{amsmath,amssymb}
\usepackage[T1]{fontenc}
\usepackage[utf8]{inputenc}
\usepackage{lmodern,microtype,fancyhdr,enumitem,hyperref,lastpage,parskip}

\hypersetup{
    colorlinks=true,
    urlcolor=black,
    linkcolor=black,
    pdftitle={CHESE11: Quality Control in Chemical Analysis -- BHU Sem I 2024-25 (NEP SEC)}
}

\pagestyle{fancy}
\fancyhf{}
\renewcommand{\headrulewidth}{0.4pt}
\renewcommand{\footrulewidth}{0.4pt}
\fancyhead[L]{\small   \textit{Banaras Hindu University}}
\fancyhead[R]{\small   \textit{CHESE11: Quality Control in Chemical Analysis (SEC)}}
\fancyfoot[C]{\small Page  \thepage\ of \pageref{LastPage}}


\newlist{parts}{enumerate}{2}
\setlist[parts,1]{label=(\alph*),leftmargin=*,itemsep=4pt,topsep=2pt}
\setlist[parts,2]{label=(\roman*),leftmargin=*,itemsep=2pt,topsep=2pt}

\newcommand{\pts}[1]{\hfill   \textit{[#1]}}


\begin{document}


\begin{center}
  {\large\bfseries BANARAS HINDU UNIVERSITY}\\[2pt]
  {
\normalsize B.Sc. (Hons) Semester I (NEP) Examination 2024-25}\\[6pt]
  \rule{0.8\linewidth}{0.5pt}\\[6pt]
  {\large\bfseries SEC --- CHEMISTRY}\\[4pt]
  {\large\bfseries Paper No. CHESE11 : Quality Control in Chemical Analysis}\\[4pt]
  {
\normalsize Time: 2:30 Hours\hfill Maximum Marks: 70}
\end{center}

\rule{\linewidth}{0.4pt}

\smallskip



\noindent (Write your Roll No. at the top immediately on the receipt of this question paper)

\smallskip



\noindent   \textit{Note: Attempt five questions including Question 1 which is compulsory. The figures in the right-hand margin indicate marks.}

\bigskip


\begin{enumerate}[label=   \textbf{\arabic*.}]

    \item Answer all of the following:
    
\begin{parts}
        \item Define and differentiate between the following:\pts{5 $ \times$ 2 = 10}
        
\begin{parts}
            \item Standard deviation and Relative standard deviation of a data set.
            \item Procedure and Protocol.
            \item Determinate and Indeterminate errors.
            \item Limit of Detection (LOD) and Limit of Quantification (LOQ).
            \item ISO and ISI.
        \end{parts}
        \item Define `Quality' and `Total Quality Management' (TQM).\pts{2}
        \item Comment on the economic benefits of using `standards' in analytical laboratories.\pts{2}
    \end{parts}

    \bigskip

    \item Answer the following:
    
\begin{parts}
        \item Following are the results of iron content analysis of a mineral: 14.50, 14.55, 14.26, 14.56, 14.53, 14.54. Calculate the following parameters for the above analysis:\pts{10}
        
\begin{parts}
            \item Standard deviation
            \item Relative standard deviation
            \item Variance
            \item Coefficient of variation
        \end{parts}
        \item How can determinate errors be minimized in analytical techniques for accurate analysis?\pts{4}
    \end{parts}

    \bigskip

    \item Answer the following:
    
\begin{parts}
        \item State any two principles of quality management.\pts{2}
        \item What is the `Six Sigma' approach to quality management? Discuss its five steps (DMAIC).\pts{6}
        \item Why is it necessary to implement a quality management system for a pharmaceutical manufacturing unit? Explain.\pts{6}
    \end{parts}

    \bigskip

    \item Answer the following:
    
\begin{parts}
        \item What does `Kaizen' mean in the context of quality management? State the basic principle on which it works.\pts{4}
        \item Why are ISO 9000 and ISO 14000 families known as `generic management system standards'? Highlight the key differences between ISO 9000 and ISO 14000.\pts{6}
        \item What is the criterion for rejection of an observation/datum in a set of analytical data?\pts{4}
    \end{parts}

    \bigskip

    \item Answer the following:
    
\begin{parts}
        \item Explain the analytical validation of results for any analytical estimation method.\pts{6}
        \item What is a normal error curve? Draw and give its mathematical expression, significance, and show the standard deviation ($\sigma$) limits in the curve.\pts{8}
    \end{parts}

    \bigskip

    \item Answer the following:
    
\begin{parts}
        \item Explain ICH guidelines for the approval of any pharmaceutical drug for human use.\pts{5}
        \item Elaborate on the role and responsibilities of `Quality Units' in a pharmaceutical manufacturing unit as per ICH guidelines.\pts{4}
        \item Give the key features of `Good Manufacturing Practices' (GMP) adapted for the production of a pharmaceutical drug.\pts{5}
    \end{parts}

    \bigskip

    \item Answer the following:
    
\begin{parts}
        \item What is `sampling'?\pts{3}
        \item How will you classify samples: (i) as per size, and (ii) as per constitution of analyte?\pts{5}
        \item How does a `laboratory' sample differ from a `real' sample?\pts{3}
        \item Write the various ways of expressing `accuracy' and `precision' of a set of analytical data.\pts{3}
    \end{parts}

\end{enumerate}

\end{document}
